\documentclass[12pt]{book}
\usepackage{precalc1}
\setcounter{chapter}{5}
\begin{document}

\chapter{Exponential Functions}

A line is built from repeated addition. Each step in the input adds the same amount to
the output. An exponential function is built from repeated multiplication. Each step in
the input multiplies the output by the same factor. That simple change from adding to
multiplying produces a very different kind of growth.

This week is about recognizing constant-ratio behavior. Populations, interest,
radioactive decay, and many physical processes are not governed by a fixed amount of
change per unit time. Instead, the amount of change depends on how much is already
present. Exponential functions are the basic language for that situation.

\[
f(x)=a\,b^{x}
\]

Here $a$ is the starting value and $b$ is the fixed multiplier. If $b>1$, repeated
multiplication produces growth. If $0<b<1$, repeated multiplication produces decay.
The rest of the chapter develops this idea, first in tables and graphs, then through the
natural base $e$, and finally through half-life.

\section{Constant Ratio, Not Constant Difference}

The first distinction is the most important one. A linear function changes by a constant
difference: each unit step adds the same amount. An exponential function changes by a
constant ratio: each unit step multiplies by the same amount. When reading data, this is
the main diagnostic. Equal differences suggest a line; equal ratios suggest an
exponential.

\begin{definition}{Exponential function}
An \emph{exponential function} has the form $f(x)=a\,b^{x}$ with base $b>0$,
$b\neq1$, and $a\neq0$. The value $a=f(0)$ is the output at $x=0$; the base $b$ is the
factor by which the output multiplies for each unit increase in $x$.
\end{definition}

\newpage
\begin{keyidea}
Linear growth adds a fixed amount per step; exponential growth multiplies by a fixed
factor per step. Over equal input steps, a linear function has equal differences and
an exponential function has equal ratios.
\end{keyidea}

\begin{example}[spotting the pattern in a table]
The outputs $3, 6, 12, 24, 48$ at $x=0,1,2,3,4$ multiply by $2$ each step, so they
fit $f(x)=3\cdot2^{x}$. The outputs $3,5,7,9,11$ instead add $2$ each step, a linear
pattern $f(x)=2x+3$. Reading whether successive outputs share a difference or a ratio
identifies the model.
\end{example}

The table test is especially useful because graphs with scatter can be visually
ambiguous. A short stretch of exponential growth may look almost linear. The ratio test
looks at the mechanism: are we adding the same amount, or multiplying by the same
factor?

\begin{center}
\begin{tikzpicture}[scale=0.78]
  \draw[->] (-2.6,0)--(2.8,0) node[right]{\scriptsize $x$};
  \draw[->] (0,-0.4)--(0,5) node[above]{\scriptsize $y$};
  \draw[pcblue,thick,domain=-2.4:2.25,smooth,samples=80] plot (\x,{2^\x});
  \node[pcblue] at (1.9,4.6) {\scriptsize $b>1$: growth};
  \draw[pcred,thick,domain=-2.25:2.4,smooth,samples=80] plot (\x,{2^(-\x)});
  \node[pcred] at (-1.6,4.6) {\scriptsize $0<b<1$: decay};
  \draw[gray,dashed] (-2.6,0)--(2.8,0);
  \fill (0,1) circle (1.6pt) node[above right]{\scriptsize $(0,1)$};
\end{tikzpicture}
\end{center}

Every exponential $b^{x}$ passes through $(0,1)$, since $b^{0}=1$, and stays strictly
positive, since a positive base raised to any power is positive. The $x$-axis is a
horizontal asymptote: the graph approaches it on one side without reaching it. With a
coefficient $a$, the point $(0,1)$ becomes $(0,a)$.

\newpage
\section{Negative Exponents and Decay}

Decay is not a new kind of formula. It is growth seen in reverse. A negative exponent is
a reciprocal:
\[
b^{-x}=\frac{1}{b^{x}}=\left(\frac1b\right)^{x}.
\]
So raising a base to negative powers is the same as raising its reciprocal to positive
powers. A base larger than $1$ run with negative exponents shrinks; equivalently, a base
between $0$ and $1$ shrinks with positive exponents. The two descriptions are the same
function.

\begin{keyidea}
$b^{-x}=\left(\dfrac1b\right)^{x}$. A negative exponent reflects the graph across the
$y$-axis and converts growth into decay. A decay function can be written with a base
above $1$ and a negative exponent, or with a base below $1$ and a positive exponent;
both name the same curve.
\end{keyidea}

\begin{example}[two ways to write decay]
The function $f(x)=2^{-x}$ halves at each step: $1,\tfrac12,\tfrac14,\tfrac18,\dots$
The same outputs come from
\[
g(x)=\left(\tfrac12\right)^{x}.
\]
Writing the decay with base $\tfrac12$ and positive exponent, or base $2$ and negative
exponent, is a matter of preference; the graph is identical.
\end{example}

The language changes, but the multiplier does not. Saying ``multiply by $1/2$ each
step'' and saying ``divide by $2$ each step'' describe the same process.

\section{The Number $e$}

One base stands out, and it arises from asking what happens when growth is compounded
more and more frequently. Imagine investing one dollar for one year at $100\%$ annual
interest. Compounded once, the year-end value is $(1+1)=2$ dollars. Compounded twice,
each half-year adds $50\%$, so the year-end factor is $(1+\tfrac12)^{2}=2.25$.
Compounded $n$ times, the year-end factor is
\[
\left(1+\frac1n\right)^{n}.
\]
As the compounding is split into ever finer periods, this factor does not run away to
infinity, nor does it stall. It climbs and settles toward a fixed limiting value.

\begin{center}
\renewcommand{\arraystretch}{1.25}
\begin{tabular}{@{}rl@{}}
\toprule
compounding $n$ & year-end factor $\left(1+\tfrac1n\right)^{n}$\\
\midrule
$1$ (annually)      & $2.000000$\\
$2$ (semiannually)  & $2.250000$\\
$4$ (quarterly)     & $2.441406$\\
$12$ (monthly)      & $2.613035$\\
$52$ (weekly)       & $2.692597$\\
$365$ (daily)       & $2.714567$\\
$8760$ (hourly)     & $2.718127$\\
$1{,}000{,}000$     & $2.718280$\\
\bottomrule
\end{tabular}
\end{center}

The values press upward against a ceiling near $2.71828$, never crossing it however
fine the compounding. That limiting value is the number $e$, an irrational constant
\[
e=2.718281828\ldots,
\]
the value the year-end factor approaches under continuous compounding.

\begin{keyidea}
The number $e$ is the limit the compounding factor $\left(1+\tfrac1n\right)^{n}$ rises
toward as the compounding is made continuous. It is the natural base for exponential
growth and decay, written $f(x)=a\,e^{kx}$, where $k>0$ gives growth and $k<0$ gives
decay.
\end{keyidea}

\begin{visual}
Splitting the year into finer compounding periods adds a little more growth each time,
but each addition is smaller than the last, and the running total is trapped beneath a
fixed height. That height is $e$. The continuous-growth curve $e^{x}$ is the smooth
result of compounding at every instant.
\end{visual}

\begin{example}[continuous growth]
A culture growing continuously at rate $k=0.5$ per hour from an initial $200$ cells is
modeled by $P(t)=200\,e^{0.5t}$. After $3$ hours,
\[
P(3)=200\,e^{1.5}\approx 200(4.4817)\approx 896
\]
cells. The base $e$ encodes growth happening at every instant rather than in discrete
steps.
\end{example}

\section{Change of Base}

Different bases can describe the same exponential curve. Changing the base does not
change the phenomenon; it changes the unit in which the rate is measured. A function
written as $3^x$ uses a multiplier of $3$ per unit step. A function written as $e^{kx}$
uses a continuous rate $k$. These are two coordinate systems for exponential behavior.

Any exponential can be rewritten in any other base, because every positive base is a
power of every other. In particular, every exponential can be written with base $e$,
which is why $e$ serves as a universal base. Writing $b=e^{k}$ for the exponent $k$ that
produces $b$, the conversion follows from the exponent rules.

\begin{keyidea}
\textbf{Change of base into $e$.} For a positive base $b$, let $k$ be the exponent
with $e^{k}=b$. Then
\[
b^{x}=\left(e^{k}\right)^{x}=e^{kx},
\]
so every exponential is a base-$e$ exponential with a rescaled exponent. The value
$k$ is denoted $\ln b$, the natural logarithm of $b$, studied next; here it is simply
the exponent that turns $e$ into $b$.
\end{keyidea}

\begin{example}[rewriting in base $e$]
Write $f(x)=3^{x}$ in base $e$. The exponent producing $3$ is $k=\ln 3\approx1.0986$
(so $e^{1.0986}\approx3$), giving
\[
3^{x}=\left(e^{\ln 3}\right)^{x}=e^{(\ln 3)x}\approx e^{1.0986\,x}.
\]
The growth base $3$ becomes the continuous rate $k\approx1.0986$.
\end{example}

\begin{example}[a decay rate in base $e$]
A substance retaining $90\%$ of its mass per year obeys $f(t)=(0.9)^{t}$. The exponent
producing $0.9$ is $k=\ln 0.9\approx-0.1054$, so
\[
f(t)\approx e^{-0.1054\,t}.
\]
The negative continuous rate confirms decay. The sign is not decoration: positive $k$
means growth, and negative $k$ means decay.
\end{example}

\newpage
\section{Half-Life and Radioactive Decay}

A radioactive element loses a fixed \emph{fraction} of its mass in each fixed span of
time. That is constant-ratio behavior, so its mass is an exponential function of time.
The natural span to measure is the \emph{half-life}: the time for the mass to fall to
half its value. After one half-life, half remains; after two, a quarter; after three, an
eighth.

\begin{definition}{Half-life}
The \emph{half-life} $T$ of a decaying quantity is the time required for it to diminish
to one-half of its current amount. Starting from $M_0$, the amount after time $t$ is
\[
M(t)=M_0\left(\tfrac12\right)^{t/T},
\]
since $t/T$ counts how many half-lives have elapsed.
\end{definition}

\begin{keyidea}
Half-life decay is constant-ratio decay with ratio $\tfrac12$ per half-life. The
exponent $t/T$ is the number of half-lives elapsed; each whole number of half-lives
multiplies the remaining amount by another factor of $\tfrac12$.
\end{keyidea}

\begin{example}[carbon-14 dating]
Carbon-14 has a half-life of about $5730$ years. A sample starting with $M_0$ grams has
\[
M(t)=M_0\left(\tfrac12\right)^{t/5730}.
\]
After $11{,}460$ years, two half-lives have passed, so
\[
M=M_0\left(\tfrac12\right)^{2}=\tfrac14 M_0.
\]
If a sample is measured at $25\%$ of its original carbon-14, it has aged two
half-lives, about $11{,}460$ years.
\end{example}

\begin{center}
\begin{tikzpicture}[xscale=0.95,yscale=1.7]
  \draw[->] (-0.3,0)--(4.6,0) node[right]{\scriptsize $t/T$ (half-lives)};
  \draw[->] (0,-0.15)--(0,1.3) node[above]{\scriptsize fraction remaining};
  \draw[pcblue,thick,domain=0:4.3,smooth,samples=80] plot (\x,{2^(-\x)});
  \foreach \n/\v in {0/1,1/0.5,2/0.25,3/0.125,4/0.0625}{
    \fill[pcred] (\n,\v) circle (1.4pt);
  }
  \node[above right] at (0,1) {\scriptsize $1$};
  \node[above right] at (1,0.5) {\scriptsize $\tfrac12$};
  \node[above right] at (2,0.25) {\scriptsize $\tfrac14$};
  \node[above right] at (3,0.125) {\scriptsize $\tfrac18$};
  \draw[gray,dashed] (-0.3,0)--(4.6,0);
\end{tikzpicture}
\end{center}

The exponent $t/T$ is a count. If $t=18$ hours and $T=6$ hours, then $t/T=3$: three
half-lives have passed. This interpretation is often easier than treating the exponent
as a mysterious algebraic decoration.

\begin{example}[solving for elapsed time without logarithms]
A medical isotope has a half-life of $6$ hours. A dose of $80$ milligrams is given; how
much remains after $18$ hours? Eighteen hours is $18/6=3$ half-lives, so
\[
M=80\left(\tfrac12\right)^{3}=80\cdot\tfrac18=10\text{ mg}.
\]
When the elapsed time is a whole number of half-lives, the answer needs no logarithms:
just halve repeatedly.
\end{example}

\begin{pitfall}
Half-life decay multiplies by $\tfrac12$ each half-life; it does not subtract a fixed
amount. After one half-life $80$ mg becomes $40$, after another it becomes $20$, not
$0$. Treating decay as repeated subtraction, rather than repeated halving, is a common
error and predicts the substance vanishing far too soon.
\end{pitfall}

\demo{In the Function Fit Lab, set the parent function to an exponential and tune
$A\,f(b(x-h))+k$ to thread a noisy point cloud. A positive horizontal scale grows and
a negative one decays; matching the curve to data is the practical form of choosing a
growth or decay rate. Compare how altering the base steepens or flattens the
climb.}{https://bobovski66.github.io/Precalculus/demos/function\_fit.html}

\section{Exercises}

\begin{exercises}
\item Decide whether each table is linear or exponential, and give a formula:
(a) outputs $5,10,20,40$ at $x=0,1,2,3$; (b) outputs $5,8,11,14$ at the same inputs;
(c) outputs $100,20,4$ at $x=0,1,2$.

\item Sketch $f(x)=3^{x}$ and $g(x)=3^{-x}$ on the same axes, mark the shared point
$(0,1)$, and identify the horizontal asymptote of each.

\item Rewrite each using a positive base and positive exponent, then state whether it
represents growth or decay: (a) $5^{-x}$; (b) $\left(\tfrac14\right)^{x}$;
(c) $2^{-3x}$.

\item Using the compounding table, explain in your own words why
$\left(1+\tfrac1n\right)^{n}$ does not grow without bound as the compounding is
refined, and state the value it approaches.

\item A bank account of \$1000 grows continuously at $4\%$ per year, modeled by
$A(t)=1000\,e^{0.04t}$. Find the balance after $10$ years.

\item Rewrite $f(x)=4^{x}$ in base $e$, giving the continuous rate $k=\ln 4$ to four
decimals.

\item Rewrite $f(t)=(0.8)^{t}$ in base $e$ and state whether the continuous rate is
positive or negative, with its meaning.

\item A radioactive sample has a half-life of $10$ years. Starting from $64$ grams,
find the amount remaining after $10$, $20$, and $30$ years without using logarithms.

\item Write the decay model $M(t)=M_0\left(\tfrac12\right)^{t/T}$ for an element with
half-life $T=1600$ years, and find the fraction remaining after $4800$ years.

\item A sample is measured at one-eighth of its original radioactive content.
Determine how many half-lives have elapsed, and if the half-life is $30$ years, the
sample's age.

\item Explain why a radioactive element's mass is an exponential function of time,
referring to the constant-ratio property and the meaning of half-life.

\item Two isotopes start with equal mass; one has half-life $2$ years, the other
$4$ years. Compute the fraction of each remaining after $8$ years, and explain which
has decayed further and why.
\end{exercises}

\end{document}
